\section{Transversal Trees  and Tree Search} \label{sec:transversal-tree}

In this section, we present an algorithm called $\ts$ 
for solving $\enumz{k}$, i.e., 
to enumerate all minimum-size transversals 
of a satisfiable $k$-CNF $F$. Let  $t=\trn(F)$ be the transversal number of $F$.
Our algorithm considers a tree,
called \emph{transversal tree}, of depth $t$
where each minimum-size transversal corresponds to at least
one leaf node at depth $t$.
Our algorithm $\ts$ traverses the tree to enumerate 
the leaves at depth $t$ corresponding to 
minimum-size transversals. 
However, 
a leaf at depth $t$ 
need not correspond to a minimum-size transversal and 
furthermore two distinct leaves at depth $t$ may 
correspond to the same
minimum-size transversal.
For these reasons, 
enumerating all leaves can take significantly more
time than the total number $\ntrn(F)$
of minimum-size transversals.
We
deal with this issue by pruning subtrees so that 
$\ts$ would only encounter 
minimum-size transversals 
that are not encountered elsewhere.
While our pruning approach is isomorphic to that of Monien-Speckenmeyer \cite{MonienS85}, 
our analysis and bound crucially depend on 
our choice of clause ordering 
and random ordering of the child nodes of the transversal tree.
In the following, we define the required concepts and present 
our constructions.

\begin{definition}[$(k,X)$-trees]
For $k\geq 1$ and  a set $X$ of variables, a $(k,X)$-tree
is a directed $k$-ary tree $T$ 
with a node and edge labeling $Q$ and a tree-edge ordering $\pi$
which satisfies the following properties.
\begin{enumerate}
\item Each edge  is directed from parent to child.
Each non-leaf node has at most $k$ children. Children of a node
are ordered from left to right according to $\pi$.
\item 
%{\bf Labels:}
Each (tree) edge $e$ is labeled with a variable $q_e\in X$.
Each node $v$ is labeled with a set $Q_v \subseteq X\cup\{\bot\}$.

\item 
%{\bf Constraints on labels:}
For each node, labels of child edges are distinct.
If $e= (u,v)$ is an edge, then 
$Q_v = Q_u\cup \{q_e\}$ or
$Q_v = Q_u\cup \{q_e, \bot\}$.
\item
Labels of  edges along any path are distinct.
\item All leaves $v$ where $\bot\notin Q_v$ 
are at the same level.
\end{enumerate}
\end{definition}
%We need the following definitions for $(k,X)$-trees.
Let $T$ be a $(k,X)$-tree with root $r$ and labeling $Q$.
For node $v$ of $T$, 
let $T_v$ denote the subtree $T_v$ of $T$ rooted at $v$.
If $u$ and $v$ are nodes of $T$ 
such that $u$ is an ancestor of $v$, $P_{uv}$ denotes the unique 
path from $u$ to $v$ in $T$.

% \begin{definition}[Stem of a subtree]
% If $v\neq r$, the \emph{stem} of $T_v$ is the tree-edge $(u,v)$ where $u$ is the parent of $v$. The \emph{stem label} of $T_v$ is the label of the tree-edge $(u,v)$.
% \end{definition}


\begin{definition}[Shoot of a tree path]
%Let $T$ be a transversal tree with labeling $Q$. 
If $u$ and $v$ are nodes of $T$ such that $u$ is an ancestor of $v$, the subgraph consisting of all the child edges of the nodes 
along the unique path from $u$ to $v$ is called the \emph{shoot} $S_{uv} = S^T_{uv}$ from $u$ to $v$. In particular $P_{uv} \subseteq S_{uv}$.
\end{definition}

For a path $P_{uv}$, the labels of the  edges along the path are called \emph{path variables} of $P_{uv}$.
For a shoot $S_{uv}$, labels of the shoot edges are called shoot variables.

\begin{definition}[Transversal tree]
Let $F=(X,\cC)$ be a $k$-CNF on the variable set $X$.
A $(k,X)$-tree $T$ rooted at $r$ with labeling $Q$ and tree-edge ordering $\pi$
is a \emph{transversal} tree for $F$ if
\begin{enumerate}
\item $Q_r = \emptyset$ if $F$ has no empty clause 
and otherwise $Q_r = \{\bot\}$, and 
\item for every node $v$, each 
minimum-size transversal of $F$ which is an 
extension of $Q_v$ will appear as 
the label of a leaf in the subtree rooted at $v$.
\end{enumerate}
\end{definition}

It is easy to see that  a transversal tree $T$ for a satisfiable
$k$-CNF $F$ has depth $\trn(F)$ and 
every leaf $v$ of $T$ such that $\bot \notin Q_v$ is at depth $\trn(F)$.
We also note that any subtree of a transversal tree is also a transversal tree.

\begin{definition}[Valid and invalid leaves]
Let $T$ be a transversal tree for a satisfiable $k$-CNF $F$.
We say that a  leaf $v$ of $T$
%at depth $\trn(F)$ 
is \emph{valid} if $Q_v$ is a minimum-size transversal of $F$.
Otherwise it is \emph{invalid}.
\end{definition}

For a transversal tree $T$ of a satisfiable $k$-CNF $F$, 
let $\trans(T)$ denote the 
collection  of minimum-size transversals associated 
with the valid leaves of $T$. The definition of transversal tree implies the following basic fact.

\begin{fact}
%Let $T$ be a transversal tree for a $k$-CNF $F$.
$\trans(F) = \trans(T)$.
\end{fact}


\subsection{Construction of Transversal Trees}

In this section, we show how to construct transversal trees for 
a satisfiable $k$-CNF 
$F=(X,\cC)$. 
The construction produces a labeling $Q$ on nodes and edges.
We will not specify a tree-edge ordering in the construction.
However, we will later select
a  left-right ordering where child nodes 
are ordered randomly and independently for each node.
The construction depends on the ordering of clauses.
Let $\Pi$ denote an ordering of the clauses in $F$.

We start with the tree $T$ with just one node $r$
(the root node) with the label  $Q_r =\emptyset$.
Assume that we are about to expand a non-leaf node $v$.
By construction, $v$ is at depth less than $\trn(F)$ and $\bot \not\in Q_v$.
Since $v$ is at depth less than $\trn(F)$, $Q_v$ is not a transversal.
%implies that $Q_v$ is not transversal) and $\bot \notin Q_v$, 
Select the first monotone clause 
$C_v = \{a_1, \ldots, a_{k'}\}$ according to the clause order $\Pi$, where $k'\le k$.
Such a monotone clause must 
exist since every clause is non-empty and since
otherwise an all-0 assignment will satisfy the formula
contradicting the fact that $Q_v$ is not a 
transversal. 
Also, it must be the case that 
$C_v\cap Q_v = \emptyset$ as none 
of the variables from $Q_v$ can appear in $C_v$.
For each $a\in C_v$,
\begin{enumerate}
\item Create a child node $v_a$ for $v$ and label the edge $(v, v_a)$ by $a$. 
\item Simplify the clauses 
by setting the variables
along the path $P_{rv_a}$ to 1. 
If there is an empty clause, set $Q_{v_a} = Q_v\cup\{a,\bot\}$ 
and $v_a$ 
will not be expanded and thus will become a leaf node.
Otherwise,
label $v_a$ by $Q_{v_a} = Q_v\cup \{a\}$. 
\item If the level
of the node is $\trn(F)$, make it a leaf node.
\item Order the child nodes of $v$ left-right according to a 
tree-edge ordering.
\end{enumerate}

\begin{proposition}
\label{prop:general-tree-existence}
The tree $T$ with the labeling as described above is a transversal tree for $F$.
\end{proposition}

\begin{proof}
Indeed each non-leaf node has at most $k$ children.
All leaves $v$ with $\bot\notin Q_v$ must be at the same level $\trn(F)$
since any node $v$ at a level smaller than $\trn(F)$ 
and does not contain $\bot$ in $Q_v$ can be expanded
and since the construction stops at level $\trn(F)$.
It is easy to verify that the labeling $Q$ has the requisite properties.
For every $v$ with $\bot \notin Q_v$
and every extension $Y$ of $Q_v$ to a minimum-size transversal, there exists a
leaf with label $Y$ in the subtree $T_v$ since there is a child edge $(v,v')$ of $v$ with label $a$ for some $a \in Y \cap C_v \ne \emptyset$ where $C_v$ is the clause used to expand $v$. Inductively we can construct a path from $v'$ to a leaf which ultimately has $Y$ as the label.
\end{proof}


\subsection{$\ts$: An Algorithm for Enumerating the Valid Leaves of $T$}

Our goal is to search the transversal tree to enumerate 
all minimum-size transversals. 
However, visiting 
all leaf nodes
may take at least $k^{\trn(F)}$ time. 
To improve the efficiency of the search,
we prune the tree during our search while
guaranteeing the enumeration of  
each minimum-size transversal exactly once.

%\begin{definition}
%A search of a $(k,X)$-tree tree $T$ is 
%\emph{non-redundant}
%if 
%$Q_u\neq Q_v$
%for every pair of 
%leaves $u\neq v$ encountered during the search.
%\end{definition}

%We define our search procedure $\ts$ for transversal trees.
Let $F$ be satisfiable $k$-CNF.
Let $\Pi$ be a clause ordering for $F$. Let $T$ be transversal
tree for $F$ constructed using $\Pi$ and some tree-edge ordering $\pi$. 
We note that for any tree-edge ordering $\pi$,
the edges of any 
shoot $S_{uv}$ (where $u$ is any ancestor of $v$)
are situated in one of three ways with respect to
the tree path from $u$ to $v$:
%\begin{itemize}
1) to the left of the tree path, 
2) to the right of the tree path, or
3) along the tree path.
Our key insight is that for any node $v$
we can determine 
whether the  subtree $T_v$ potentially contains 
any \emph{new }
minimum-size transversals by considering the labels of the edges
in the shoot $S_{rv}$.


Our $\ts$ starts with the root node $r$ of $T$.
Assume that we are currently visiting the node $v$.
Let $T_v$ be a subtree of $T$ rooted at $v$.
If $v$ is not a leaf, let $C_v$ be the monotone clause used for expanding
the node $v$ (based on the clause order $\Pi$) and   $a_1, \ldots, a_{k'}$ be the ordering of its variables
according to $\pi$ for some $k' \le k$. For $1\leq i \leq k'$, 
let $v_{a_i}$ be the $i$-th child node of $v$. The edge $e_i=(v,v_{a_i})$ is labeled with $a_i$.
The search procedure $\ts$ starting at node $v$ works as follows:

\begin{itemize}
\item  If $v$ is a leaf,
output $Q_v$ if it is a transversal.
In any case, return to the parent.
\item Otherwise, process the children of $v$ in order. 
Let $T_{i} = T_{v_{a_i}}$ be the transversal tree 
rooted at the child $v_{a_i}$ for $a_i \in C_v$. 
Prune the subtree $T_i$ if and only if 
the shoot $S_{rv}$ contains an edge $e'=(u',v')$ 
where $u'$ is ancestor of $v$ such that $Q_{e_i} = Q_{e'}$, and
the edge $e'$ appears to the left of the path $P_{rv}$.
Search the tree $T_i$ if it is not pruned.

\end{itemize}

\begin{fact}
For any clause ordering $\Pi$
and tree-edge ordering $\pi$, 
$\ts$ outputs all minimum-size transversals of $F$ exactly once.
\end{fact}

\noindent
\subsection{Canonical Clause Ordering and Random Tree Edge Ordering}

The time complexity of $\ts$ is bounded by the number of leaf nodes it visits.
While we know that $\ts$ outputs 
all minimum-size transversals without redundancy, 
it is much less clear how to analyze its complexity.
We need two ideas to analyze $\ts$ to get a good bound.
The first idea is a \emph{canonical} clause ordering  $\Pi$
in which a sequence of  maximally disjoint
monotone clauses will precede all other clauses.
The second idea is a random $\pi$, that is, 
a tree-edge ordering that orders 
the children of every node in the transversal tree
uniformly 
and independently at random.

%In the following two sections, we analyze our algorithm $\ts$ on 
%random and canonical transversal trees and present our bounds.
